\documentclass{article}
\usepackage[utf8]{inputenc}
\usepackage{graphicx} % Required for inserting images
\usepackage[english,russian]{babel}
\usepackage{hyperref}
\usepackage{amssymb}

\begin{document}

\section*{Метрические методы. Метрики качества \newline 10 баллов. +2 бонусных балла}

\subsection*{Задача 1. (1 балл)}

    Оценим время работы алгоритма ближайших соседей по количеству операций. Пусть $X$~--- обучающая выборка размера $n$, $Y$~--- тестовая выборка размера $m$. Размерность признакового пространства $d$. 
    
    \noindent
    Таким образом, $X \in \mathbb{R}^{n \times d}$, а $Y \in \mathbb{R}^{m \times d}$.
    
    \noindent
    Квадрат евклидова расстояния между объектами $x_i$ и $y_j$ записывается как:
    $$
        \rho(x_i, y_j) = \sum_{k=1}^d (x_i^k - y_j^k)^2.        
    $$
    \begin{itemize}
        \item Определите количество операций, необходимое для подсчета всех попарных растояний в наивном случае.
        
        \item Предложите способ, с помощью которого можно уменьшить количество операций. Оцените количество операций для предложенного метода.
    \end{itemize}

\noindent
\textbf{Решение}:

    \begin{itemize}
        \item На вычисление растояния двух точек мы тратим $3d$ операции. Так как мы считаем расстояния для $n * m$ пар точек то в итоге происходит $3nmd$
        
        \item Заметим что $\| x_i - y_j\|^2 = \|x_i\|^2 + \|y_i \|^2 -2 \langle x_i, y_i \rangle $
        
        в таком случае посчитаем заранее нормы всех векторов $X$. Норма $y_i$ считается за $2d$ операции (возведения в квадраты и сумма), таких векторов $m$. скалярное произведение считается за $2d$ операции (перемножение и сумма), таких пар всего $n*m$
        итого операции = $2md + 2nmd$
        
    \end{itemize}

\subsection*{Задача 2. (2 балла)}

    Дано $n$ объектов, распределённых равномерно внутри $d$-мерного единичного шара с центром в нуле.
    \begin{itemize}
        \item Найдите выражение для медианы расстояния от начала координат до ближайшего объекта.
        \item Проинтерпретируйте полученный результат в терминах применимости метода ближайшего соседа в различных ситуациях.
    \end{itemize}
    
    \noindent
    Считайте, что метрика в задаче евклидова.
    
    \noindent
    \emph{Указание: попробуйте смоделировать событие и посчитать его вероятность в терминах функций распределения.}

\noindent
\textbf{Решение}:
    \begin{itemize}
        \item Рассмотрим расспределение такой случайной величины: $min( \{ \| x_i \| \}_{i=0} )$ где $x_i$ точки d-мерного шара.

$$ F(x) = P(min( \{ \| x_i \| \}_{i=0} ) \leq x) = $$
$$ 1 - P(min(\{ \| x_i \| \}_{i=0} > x) $$
$$ 1 - P( \| x_i \| > x)^n $$
$$ 1 - (1 - P( \| x_i \| \leq x))^n $$

Заметим что $P( \| x_i \| \leq x)$ вероятность попадания точки в шар радиусом $x$. обьем шара $V_d(r) = c_d r^d$ где $c_d = \frac{\pi^\frac{d}{2}}{\Gamma(\frac{d}{2}) + 1}$

$P( \| x_i \| \leq x) = \frac{V_d(r)}{V_d(1)} = r^d$

Получаем

$$F(x) = 1- (1 - x^d)^n$$

$$F(x_{med}) = \frac{1}{2}$$

$$x_{med} = (1 - 2^{-\frac{1}{n}})^{\frac{1}{d}}$$
        \item при росте размерности, медиана минимального расстояния будет стемится к 1. что значит что все почти все наши обьекты будут находится на границе шара. Становится труднее понимать что значит "близко".

        
    \end{itemize}


\subsection*{Задача 3. (3 балла)}

    Решается задача классификации с помощью алгоритма ближайших соседей (метрика евклидова). Для тестового объекта $z$ ближайшим соседом с расстоянием $\rho_x$ является $x$, вторым ближайшим соседом с расстоянием $\rho_y$ является объект $y$. Остальные объекты обучающей выборки находятся от $z$ на достаточно большом расстоянии.

    Ко всем объектам добавляется новый признак: для $z$ и $y$ значение признака распределёно равномерно на отрезке $[-1, 1]$, для всех остальных объектов значение признака равно нулю. 
    \begin{itemize}
        \item Посчитайте вероятность того, что теперь ближайшим соседом для $z$ будет не $x$, а $y$.
        \item Проинтерпретируйте полученный результат в терминах применимости метода ближайших соседей.
    \end{itemize}

    \noindent
    \emph{Указание: возможно в этой задаче пригодится знание криволинейных интегралов.}

\noindent
\textbf{Решение}:

Исходные квадраты расстояний:

До объекта $x$: $\rho_x^2$

До объекта $y$: $\rho_y^2$, причём $\rho_x < \rho_y$

После добавления нового признака:

Для $z$ и $y$: значение признака равномерно на $[-1, 1]$

Для $x$: значение признака равно 0


пусть $u$ — значение признака для $z$, $v$ — значение признака для $y$


Новые квадраты расстояний:
$$
\rho_{\text{new}}(z, x)^2 = \rho_x^2 + u^2
$$
$$
\rho_{\text{new}}(z, y)^2 = \rho_y^2 + (u - v)^2
$$

Условие того, что $y$ становится ближайшим соседом:
$$
\rho_{\text{new}}(z, y) < \rho_{\text{new}}(z, x)
$$

Подставляем выражения:
$$
\rho_y^2 + (u - v)^2 < \rho_x^2 + u^2
$$

Упрощаем:
$$
\rho_y^2 + u^2 - 2uv + v^2 < \rho_x^2 + u^2 
$$
$$
-2uv + v^2 < \rho_x^2 - \rho_y^2 
$$
$$
2uv > v^2 + (\rho_y^2 - \rho_x^2)
$$

Вводим обозначение $\delta = \rho_y^2 - \rho_x^2 > 0$:
$$
2uv > v^2 + \delta
$$

Признаки $u$ и $v$ независимы и равномерно распределены на $[-1, 1]$. Их совместная плотность вероятности:
$$
f(u, v) = \frac{1}{4}, \quad (u, v) \in [-1, 1] \times [-1, 1]
$$

Вероятность того, что $y$ станет ближайшим соседом:
$$
P = P(2uv > v^2 + \delta) = \frac{1}{4} \iint\limits_{2uv > v^2 + \delta} du\,dv
$$

Оставим выкладки 
$$
P = \frac{I_1}{2} = \frac{1}{8} + \frac{\sqrt{1 - \delta}}{4} - \frac{\delta}{8} + \frac{\delta}{4} \ln(1 - \sqrt{1 - \delta})
$$

Полученный результат показывает уязвимость метода к добавлению ложных (шумовых) признаков:
При $\delta \to 0$ (исходные расстояния почти равны) вероятность $P \to \frac{3}{8}$. Это означает, что даже небольшой шумовыой признак может с не малой вероятностью изменить ближайшего соседа.

\subsection*{Задача 4. (1 балл)}

    Докажите, что ROC-AUC случайного классификатора равен $0.5$.

\noindent
\textbf{Решение}

по определению TPR =$ \frac{TP}{TP + FN}$ = доля верно предсказанных позитивных классов. Ее можно переписать вот так:

\noindent
$TPR(x) = P(f(x, w, x_0) > x | x = True)$

где f - случайный классификатор из условия. Аналогично

\noindent
$FPR(x) = P(f(x, w, w_0) > x | x = False)$

так как f - классифицирует случайно, то 

\noindent
$P(f(x, w, x_0) > x | x = True) = P(f(x, w, x_0) > x)$

\noindent
$P(f(x, w, x_0) > x | x = False) = P(f(x, w, x_0) > x)$

То есть $TPR(x) = FPR(x)$

$$ROC-AUC =$$
$$= \int_1^0 TPR(x) FPR'(x)dx =$$
$$= \int_1^0 TPR d(FPR) = $$
$$= \int_1^0 TPR d(TPR) = \frac{1}{2}$$

Что и требовалось доказать

\subsection*{Задача 5. (2 балла)}

    Пусть, $a = a(x)$ ответ алгоритма. На сколько может уменьшиться ROC-AUC при использовании функции $\min(a, 0.5)$ над оценками алгоритма?

\noindent
\textbf{Решение}

Посмотрим TPR FPR нового алгоритма a'

$$TPR_{new}(x) = TPR(x),\quad \text{if}  \quad x \leq 0.5$$

$$TPR_{new}(x) = 0,\quad \text{if}  \quad x > 0.5$$

$$FPR_{new}(x) = FPR(x),\quad \text{if}  \quad x \leq 0.5$$

$$FPR_{new}(x) = 0,\quad \text{if}  \quad x > 0.5$$

Тогда ROC' будет состоять из отрезка соединяющего (0,0) (TPR(0.5), FPR(0.5)) и кривой совпадающей c ROC на $t \in [0, 0.5]$

и по сути потерянная часть ROC-AUC:

$$ROC-AUC - ROC-AUC_{new}= \int_1^{0.5} TPR(x) FPR'(x)dx - \frac{1}{8}$$

\noindent
\textbf{Замечание}

На самом деле, такое преобразование алгоритма опускает начальную часть кривой до прямой TPR = FPR что для хорошего алгоритма может быть критичным, ведь на высоких порогах, он мог работать хорошо, и соответсвенно его ROC кривая выглядила бы более впуклой. Но после изменения алгоритма, мы бы потеряли то самое преимущество на высоких порогах


\subsection*{Задача 6. (3 балла)}

    Подробнее ознакомьтесь с материалом по ROC-AUC по \href{https://alexanderdyakonov.wordpress.com/2017/07/28/auc-roc-%D0%BF%D0%BB%D0%BE%D1%89%D0%B0%D0%B4%D1%8C-%D0%BF%D0%BE%D0%B4-%D0%BA%D1%80%D0%B8%D0%B2%D0%BE%D0%B9-%D0%BE%D1%88%D0%B8%D0%B1%D0%BE%D0%BA/}{ссылке} и решите следующую задачу:

    \noindent
    Пусть на ответах алгоритма $m$ (принимающих значения от $0$ до $1$) задано распределение объектов класса $1$ (доля объектов класса $1$ в зависимости от ответа алгоритма) следующим образом:
    $$
        \mathbb{P} (m \in [a, b] \ | \ y = 1) = \int_a^b p(z)dz.
    $$

    \noindent
    Распределение объектов класса 0 задаётся так:
    $$
        \mathbb{P} (m \in [a, b] \ | \ y = 0) = \int_a^b (2 - p(z))dz,
    $$
    где $p(z) = -1.5z^2 + 3z$.

    \noindent
    Найдите вероятностные оценки на величины TPR, FPR и ROC-AUC.


\noindent
\textbf{Решение}

\noindent TPR(t) = $P(m \geq t | y = 1) = \int_t^1 p(z)dz$

\noindent FPR(t) = $P(m \geq t | y = 0) = \int_t^1 (2 - p(z))dz$

\noindent Подставляем $p(z) = -1.5z^2 + 3z$:

$$TPR(t) = \int_t^1 (-1.5z^2 + 3z)dz = 0.5t^3 - 1.5t^2 + 1$$

$$FPR(t) = \int_t^1 (1.5z^2 - 3z + 2)dz = -0.5t^3 + 1.5t^2 - 2t + 1$$

\noindent
Найдём ROC-AUC:


$$ROC-AUC = \int_1^0 TPR(t) FPR'(t) dt$$

$$ROC-AUC= \int_0^1 (0.5t^3 - 1.5t^2 + 1)(1.5t^2 - 3t + 2)dt =  0.125 - 0.75 + 1.375 - 0.5 - 1.5 + 2 = 0.75$$
\end{document}
